\section{Open questions raised (and not answered) by this replication}

\begin{enumerate}

\item \textbf{Open-shell extension.} Does CS-VQE extend cleanly to
open-shell molecules (spin doublets/triplets, e.g.\ OH, O$_2$, CH$_3$)
where the Hartree--Fock reference is no longer a single closed-shell
determinant, and does the noncontextual partition still capture the
dominant static correlation? \emph{Basis:} the paper and this
replication only test closed-shell singlets in minimal basis; the
noncontextual quasi-quantized model's accuracy for radicals is
unproven. \emph{Next steps:} build STO-3G JW Hamiltonian for OH (7
electrons, 12 qubits) with ROHF/UHF; enumerate noncontextual
partitions and record min $q$ for 1 mHa vs ROHF-CISD; compare CS-VQE
reduction ratio for OH vs the paper's H$_2$O row; repeat for O$_2$
triplet as a stronger test on a degenerate ground manifold.

\item \textbf{Composition with tapering / symmetry reduction.} How
does CS-VQE compose with Z$_2$ symmetry tapering (Bravyi et al.\ 2017)
and point-group symmetry reduction --- are reductions additive or does
the noncontextual set already absorb the same symmetries?
\emph{Basis:} both techniques exploit stabilizer subgroups; the paper
does not test composition and no downstream benchmark exists (as of
2026). \emph{Next steps:} for H$_2$O/STO-3G (14 qubits pre-taper),
compare qubit budget after tapering-alone, CS-alone (exhaustive up to
$q\!=\!6$), taper-then-CS, and CS-then-taper; formally check whether
our greedy nc-generators intersect the Z$_2$ taper generators.

\item \textbf{Strongly-correlated regime.} Does CS-VQE remain useful
for strongly-correlated systems (Cr$_2$, stretched N$_2$, TM dimers)
where the noncontextual approximation should fail because multiple
determinants have comparable weight, or does it degrade to $q\!\to\!n$
as static correlation grows? \emph{Basis:} the paper's H$_2$ at
0.7414~\AA\ is weakly correlated; the strongly-correlated regime is
the actually-interesting one, and it is not tested. \emph{Next steps:}
extend our H$_2$ script to scan 0.5--3.5~\AA\ in 0.1~\AA\ steps and
plot $q_{\text{required}}(R)$; repeat for stretched N$_2$/STO-3G;
correlate $q_{\text{required}}$ with the D1 diagnostic and
$|c_{HF}|^2$ in the FCI expansion.

\item \textbf{ML-selected noncontextual partitions.} Can a learned
model (GNN over the Pauli-anticommutation graph, or a Transformer over
term coefficients) predict a near-optimal noncontextual partition
faster than the paper's greedy heuristic, and does the ML-selected
partition generalize across molecular families? \emph{Basis:} our
exhaustive H$_2$ search found partitions the greedy heuristic misses
(the $q\!=\!0$ partition using YXXY); exhaustive costs $2^N$ and is
infeasible for $N\!>\!\sim\!20$. \emph{Next steps:} generate $\sim$10$^4$
labelled small Hamiltonians ($n\!=\!4$--$8$); train a GNN on the
anticommutation graph to predict per-term nc/c membership; on
held-out H$_2$O/NH$_3$ compare exhaustive, greedy, and GNN-selected
partitions against FCI.

\item \textbf{Noise robustness of the reduced-qubit circuit.} Is
CS-VQE more or less robust to realistic device noise (T$_1$/T$_2$,
readout, crosstalk) than full VQE at the same target accuracy, once
the smaller circuit is actually run on a parameterized ansatz --- not
just classically diagonalized as in the paper's headline sweeps and
this replication? \emph{Basis:} the entire near-term case for CS-VQE
rests on this and it is directly untested. \emph{Next steps:} build
hardware-efficient ansätze at $q\!=\!1$ (CS-VQE) and $q\!=\!4$ (full)
for H$_2$; simulate with Qiskit Aer's noise model calibrated to a
public IBM Falcon backend; run 100 VQE optimizations per (ansatz,
shots $\in\{1024, 8192, 65536\}$); extend to LiH/STO-3G to test whether
the noise-robustness advantage grows with system size.

\end{enumerate}
